\cleardoublepage
\chapter{Opis wybranych algorytmów sortowania równoległego}
\label{cha:WybraneAlgorytmy}

Dział poświęcony wybranym algorytmom zawierający opis, schemat działania, sposób dzielenia i łączenia, złożoność oraz 
potencjalne problemy.

\section{Parallel Quicksort}
\label{sec:ParallelQuicksort}

Quicksort należy do rodziny algorytmów opartych na strategii ``dziel i zwyciężaj'' (ang. \textit{divide-and-conquer})~\cite{bib:CLRS}. 
W kroku podziału wybierany jest element rozdzielający (ang. \textit{pivot}), który dzieli wejściowy ciąg $A[p \ldots r]$ na dwa podciągi. 
Po zakończeniu partycjonowania pivot trafia na pozycję $A[q]$, przy czym wszystkie elementy w podciągu $A[p \ldots q-1]$ są mniejsze 
lub równe $A[q]$, a wszystkie elementy w podciągu $A[q+1 \ldots r]$ są od niego większe. 
Następnie oba podciągi są rekurencyjnie sortowane aż do uzyskania podciągów o długości 1~\cite{bib:Božidar2015}.

Średnia złożoność obliczeniowa Quicksortu wynosi $O(n \log n)$, natomiast w przypadku pesymistycznym 
związanym z silnie niezrównoważonym podziałem może wzrosnąć do $O(n^2)$~\cite{bib:CLRS}.

W wersji równoległej algorytm wykorzystuje fakt, że po wykonaniu partycjonowania powstałe podciągi mogą być sortowane niezależnie od siebie. 
Dzięki temu możliwe jest równoległe przetwarzanie obu części tablicy przez osobne wątki. W praktycznych realizacjach poziom równoległości 
może być ograniczany poprzez ustalenie maksymalnej głębokości równoległej rekurencji, co umożliwia zmniejszenie nadmiernego narzutu związanego 
z tworzeniem i synchronizacją wątków. Po osiągnięciu założonej głębokości dalsze sortowanie odbywa się sekwencyjnie.

Jednym z głównych wyzwań wersji równoległej Quicksortu pozostaje nierównomierny podział danych, który występuje w sytuacji, 
gdy pivot jest źle dobrany. Problem ten, w połączeniu z narzutem związanym z tworzeniem i synchronizacją wątków, 
może istotnie ograniczać skalowalność algorytmu~\cite{bib:Maus}. 
Algorytm sprawdza się najlepiej przy danych zrównoważonych i dobrym wyborze pivota,
traci natomiast na efektywności m.in. gdy podziały stają się nierówne.

\begin{figure}[htbp]
  \centering \includegraphics[width=0.85\textwidth]{images/chapter3/parallelQuicksort.png}
  \caption{Schemat działania Quicksort.} \label{fig:qs-partition}
\end{figure}

\section{Parallel Merge Sort}
\label{sec:ParallelMergeSort}

Merge Sort należy do algorytmów opartych na strategii ``dziel i zwyciężaj'' (ang. \textit{divide-and-conquer})~\cite{bib:CLRS4}. 
Algorytm dzieli wejściowy ciąg na dwa podciągi o zbliżonej długości, 
następnie rekurencyjnie sortuje obie części, a na końcu scala je w jeden posortowany ciąg.

Operacja scalania stanowi kluczowy etap algorytmu i dla dwóch posortowanych podciągów 
o łącznej liczbie $n$ elementów może zostać wykonana w czasie $O(n)$.
Jeżeli zastosowana procedura scalania jest stabilna, cały algorytm Merge Sort również zachowuje stabilność~\cite{bib:Božidar2015}.

Złożoność czasowa Merge Sort wynosi $O(n \log n)$.
W przeciwieństwie do Quicksortu jej rząd nie zależy od początkowego uporządkowania danych,
ponieważ algorytm w kolejnych krokach dzieli problem na dwa podproblemy o zbliżonych rozmiarach
i wykonuje liniową operację scalania~\cite{bib:CLRS4}.

W wersji równoległej wykorzystuje się niezależność dwóch podproblemów powstałych po podziale danych.
Rekurencyjne sortowanie obu części może dzięki temu odbywać się równolegle.
Przed rozpoczęciem scalania konieczne jest jednak zakończenie sortowania obu podciągów,
ponieważ etap scalania wykorzystuje wyniki obu wcześniejszych operacji~\cite{bib:CLRS4}.
Samo scalanie może być realizowane sekwencyjnie lub równolegle.

Jednym z ograniczeń Parallel Merge Sort jest etap scalania.
Wariant wykorzystujący równoległe sortowanie podciągów, lecz sekwencyjne scalanie, oferuje ograniczony poziom równoległości, 
ponieważ sekwencyjna operacja scalania staje się wąskim gardłem algorytmu~\cite{bib:CLRS4}.
Zastosowanie równoległego scalania pozwala zwiększyć dostępny poziom równoległości,
jednak wiąże się z bardziej rozbudowanym sposobem łączenia posortowanych części.
W praktycznych realizacjach dla odpowiednio małych podproblemów można również przejść
do wykonania sekwencyjnego w celu ograniczenia narzutu związanego z równoległością~\cite{bib:CLRS4}.

\begin{figure}[htbp]
    \centering
    \includegraphics[width=0.85\textwidth]{images/chapter3/parallelMergeSort.jpg}
    \caption{Schemat działania Merge Sort.}
    \label{fig:ms-merge}
\end{figure}

\section{Parallel Bucket Sort}
\label{sec:ParallelBucketSort}

Bucket Sort jest algorytmem sortowania, którego działanie opiera się na podziale zakresu wartości danych wejściowych 
na określoną liczbę kubełków (ang. \textit{buckets}). Każdy element wejściowy jest przypisywany do odpowiedniego kubełka 
na podstawie swojej wartości, a następnie zawartość poszczególnych kubełków jest sortowana niezależnie. 
Na końcu posortowane kubełki są łączone w ustalonej kolejności, tworząc wynikowy ciąg~\cite{bib:CLRS4}.

Przy założeniu odpowiedniego rozkładu danych wejściowych Bucket Sort może osiągać średnią złożoność czasową $O(n)$. 
W klasycznej analizie zakłada się równomierny i niezależny rozkład elementów, dzięki czemu można oczekiwać, 
że liczba elementów przypadających na poszczególne kubełki będzie zbliżona. 
Każdy kubełek może być następnie sortowany za pomocą innego algorytmu, na przykład Insertion Sort, 
a ostateczny wynik powstaje przez konkatenację kubełków w kolejności~\cite{bib:CLRS4}.

W wersji równoległej możliwe jest zrównoleglenie zarówno etapu rozmieszczania elementów w kubełkach, 
jak i sortowania ich zawartości. Dane wejściowe mogą zostać podzielone na części, które są przetwarzane równolegle 
przez niezależne wątki. Po zakończeniu rozmieszczania elementów kubełki mogą zostać podzielone na grupy, 
których sortowanie również odbywa się równolegle~\cite{bib:Hong}. 
Po zakończeniu tej fazy posortowane kubełki są łączone w odpowiedniej kolejności w~jeden uporządkowany ciąg.

Efektywność Parallel Bucket Sort jest w dużym stopniu zależna od sposobu rozłożenia elementów pomiędzy kubełkami. 
Przy równomiernym wypełnieniu kubełków możliwe jest bardziej równomierne rozłożenie pracy pomiędzy jednostki wykonawcze. 
Jeżeli natomiast część kubełków zawiera znacznie więcej elementów niż pozostałe, może prowadzić to do nierównomiernego obciążenia 
i~ograniczenia uzyskiwanego przyspieszenia. Dodatkowym ograniczeniem jest narzut związany z realizacją obliczeń równoległych, 
w tym z synchronizacją oraz organizacją pracy pomiędzy jednostkami wykonawczymi. 
Dla małych zbiorów danych koszt ten może przewyższać korzyści wynikające ze zrównoleglenia, 
natomiast dla większych zbiorów możliwe jest uzyskanie lepszej wydajności~\cite{bib:Hong}.

\begin{figure}[htbp]
    \centering
    \includegraphics[width=0.85\textwidth]{images/chapter3/parallelBucketSort.png}
    \caption{Schemat działania Bucket Sort.}
    \label{fig:bs-buckets}
\end{figure}

\section{Parallel Samplesort}
\label{sec:ParallelSamplesort}

